\documentclass{article}
\usepackage{graphicx} % Required for inserting images
\usepackage[utf8]{inputenc}
\usepackage[T2A]{fontenc}
\usepackage{amsmath}
\usepackage{pgfplots}
\usepackage{longtable}
\usepackage{tabularx}
\usepackage{makecell}
\usepackage{array}
\usepackage{indentfirst}
\usepackage{icomma}
\usepackage{geometry}
\usepackage{booktabs}
\usepackage{multirow}
\usepackage{multicol}
\usepackage{float}
\usepackage{diagbox}
\usepackage{gensymb}
\geometry{a4paper, margin=1in}

\renewcommand{\figurename}{Изображение}
\renewcommand{\tablename}{Таблица}

\title{Лабораторная работа №2.07

Определение показателя адиабаты воздуха с использованием осциллятора Фламмерсфельда}
\author{Топунов Антон, Z3142}
\date{21.03.2026}

\begin{document}

\maketitle
\clearpage
\section{Цели и задачи}
\subsection{Цели}
\begin{enumerate}
    \item Определение показателя адиабаты воздуха.
    \item Определение эффективного числа степеней свободы для молекул воздуха.
\end{enumerate}
\subsection{Задачи}
\begin{enumerate}
    \item Измерение периода колебаний осциллятора Фламмерсфельда.
\end{enumerate}
\section{Теоретическое введение}
Период колебаний для одного запуска:
\begin{gather*}
    T_j = \frac{ t_j}{ N}
\end{gather*}
\begin{itemize}
    \item $N$ - количество колебаний
    \item $t_j$ - время $N$ колебаний для $j$-го запуска
\end{itemize}
Среднее значение периода колебаний и его погрешность:
\begin{gather*}
    T = \frac{ \sum\limits_{ j=1}^{ M}T_j}{ M} \\
    \Delta T_{ \text{сист}} = \frac{ \Delta{t}}{ N}\\
    \Delta T_{ \text{случ}} = t_{ \alpha, M}\sqrt{\frac{ \sum\limits_{ j=1}^{M }(T_j - T)^2}{ M(M-1)} }\\
    \Delta T = \sqrt{\Delta T_{ \text{сист}}^2 + \Delta T_{ \text{случ}}^2}
\end{gather*}
\begin{itemize}
    \item $M$ - количество запусков установки (количество серий измерений)
    \item $t_{ \alpha, M}$ - коэффициент Стьюдента
    \item $\Delta T_{ \text{сист}}$ - систематическая погрешность периода
    \item $\Delta T_{ \text{случ}}$ - погрешность разоброса
\end{itemize}

Показатель адиабаты воздуха:
\begin{gather*}
    \gamma = \frac{ 4mV}{ T^2p_0r^4}\\
    \Delta \gamma = \gamma\sqrt{\bigg( \frac{ \Delta{m}}{m } \bigg)^2 + \bigg( \frac{ \Delta{V}}{V } \bigg)^2
        +\bigg( \frac{ 2\Delta{T}}{T } \bigg)^2 + \bigg( \frac{ \Delta{p_0}}{p_0 } \bigg)^2 + \bigg( \frac{ 4\Delta{r}}{r } \bigg)^2} 
\end{gather*}
\begin{itemize}
    \item $m$ - масса колеблющегося поршня
    \item $r$ - внутренний радиус трубки
    \item $V$ - рабочий обьём колбы
    \item $p_0$ - атмосферное давление
\end{itemize}

Теоретическое значение показателя адиабаты воздуха:
\begin{gather*}
    \gamma = \frac{ i+2}{ i}
\end{gather*}
\begin{itemize}
    \item $i$ - число эффективных степеней свободы молекул
\end{itemize}

Из этого выражения и экспериментального значения $\gamma$ можно получить выражение для числа степеней свободы вохдуха:
\begin{gather*}
    i = \frac{ 2}{ \gamma-1}\\
    \Delta i = \frac{ 2\Delta{\gamma}}{ (\gamma-1)^2} 
\end{gather*}
\section{Схема установки}
В этой лабораторной работе используется осциллятор Фламмерсфельда:
\begin{figure}[H]
   \centering
   \includegraphics*[width=0.8\linewidth]{scheme1.PNG}
   \caption{Внешний вид экспериментальной установки}
\end{figure}
\begin{enumerate}
    \item осциллятор (легкий поршень)
    \item трубка
    \item колба
    \item отверстие
    \item балластный объем
    \item компрессор
    \item редукционный клапан
    \item оптические ворота
    \item секундомер
\end{enumerate}

Осциллятор (поршень) 1 движется в трубке 2, которая составляет единое целое с
колбой 3. Объем колбы является рабочим объемом $V$. В стенке 8
трубки около положения равновесия осциллятора сделано отверстие 4. Поднимаясь вверх, осциллятор открывает отверстие,
и давление в колбе становится равным атмосферному. Выше отверстия
осциллятор движется только под действием силы тяжести.
При движении вниз осциллятор перекрывает отверстие и сжимает воздух в колбе, после чего цикл повторяется.
Для поддержания постоянного давления в колбе и восполнения потерь воздуха через отверстие 4 служит балластный объем 5,
в который постоянно подкачивается воздух компрессором 6. Регулируя подачу воздуха с помощью редукционного клапана 7
можно добиться постоянства амплитуды колебаний осциллятора.
При своем движении осциллятор пересекает луч фотодатчика оптических ворот 8.
Общее время колебаний измеряется секундомером 9.
\section{Используемые приборы, инструментальные погрешности}
\begin{itemize}
    \item Секундомер $\Delta t = 0.02\text{ c}$
    \item Лабораторный манометр $\Delta p = 0.1 \text{ кПа}$
    \item Количество колебаний измеряется оптическими воротами 
\end{itemize}
\section{Данные прямых измерений, табличные значения}
\begin{table}[H]
   \centering
   \caption{Прямые измерения времений $N$ колебаний}
   \begin{tabular}{|c|c|} \hline
       $j$ & $t_{ j}$, с \\ \hline
       $1$ & $16.83$ \\ \hline
       $2$ & $16.72$ \\ \hline
       $3$ & $16.81$ \\ \hline
       $4$ & $16.75$ \\ \hline
       $5$ & $16.76$ \\ \hline
   \end{tabular}
\end{table}

Количество серий измерений:
\begin{gather*}
    M = 5
\end{gather*}

Коэффициент Стьюдента для доверительной вероятносит $\alpha = 0.95$ и количества измерений $M = 10$:
\begin{gather*}
    t_{ 0.95, 5} = 2.78
\end{gather*}

Количество колебаний за один запуск установки:
\begin{gather*}
    N = 50
\end{gather*}

Масса поршня:

\begin{gather*}
    m = (4.59\pm0.02)\text{ г}
\end{gather*}

Внутренний радиус трубки:
\begin{gather*}
    r = (5.95\pm0.02) \text{ мм}    
\end{gather*}

Рабочий объём колбы
\begin{gather*}
    V = (1.14\pm0.01) \text{ дм}^3
\end{gather*}

Атмосферное давление:
\begin{gather*}
    p_0 = (103.9\pm0.1) \text{ кПа}
\end{gather*}

\section{Расчёты и графики}
\begin{table}[H]
   \centering
   \caption{Период колебаний для каждого запуска}
   \begin{tabular}{|c|c|} \hline
        $j$ & $T$, c \\ \hline
        $1$ & $0.3366$ \\ \hline
        $2$ & $0.3344$ \\ \hline
        $3$ & $0.3362$ \\ \hline
        $4$ & $0.3350$ \\ \hline
        $5$ & $0.3352$ \\ \hline
\end{tabular}
\end{table}

Средний период колебаний и его погрешность:
\begin{gather*}
        T = \frac{ \sum\limits_{ j=1}^{ 5}T_j}{ 5} \approx  335.5\cdot10^{-3}\text{ с}\\
    \Delta T_{ \text{сист}} = \frac{ 0.02}{ 50} = 0.4\cdot10^{-3}\text{ с}\\
    \Delta T_{ \text{случ}} = 2.78\cdot\sqrt{\frac{ \sum\limits_{ j=1}^{M }(T_j - 0.3355)^2}{ 4\cdot5} }\approx1.1\cdot10^{-3}\text{ с}\\
    \Delta T = \sqrt{(0.4\cdot10^{-3})^2 + (1.1\cdot10^{-3})^2}\approx1.2\cdot10^{-3}\text{ с}
\end{gather*}

Показатель адиабаты:
\begin{gather*}
    \gamma = \frac{ 4\cdot4.59\cdot10^{-3}\cdot1.14\cdot10^{-3}}{ (335.5\cdot10^{-3})^2\cdot 103.9\cdot10^{3}\cdot(5.95\cdot10^{-3})^4}
        \approx 1.428\\
    \Delta \gamma = 1.428\cdot\sqrt{\bigg( \frac{ 0.02\cdot10^{-3}}{4.59\cdot10^{-3} } \bigg)^2 +
        \bigg( \frac{ 0.01\cdot10^{-3}}{1.14\cdot10^{-3} } \bigg)^2
        + \bigg( \frac{ 2\cdot1.2\cdot10^{-3}}{335.5\cdot10^{-3}} \bigg)^2 +
        \bigg( \frac{ 0.1\cdot10^3}{103.9\cdot10^{3} } \bigg)^2 +
        \bigg( \frac{ 4\cdot0.02\cdot10^{-3}}{5.95\cdot10^{-3} } \bigg)^2} 
        \approx0.026
\end{gather*}

Количество эффективных степеней свободы:
\begin{gather*}
    i = \frac{ 2}{ 1.428-1} \approx 4.67 \\
    \Delta i = \frac{ 2\cdot0.026}{ (1.428-1)^2} \approx 0.28 
\end{gather*}

\section{Результаты и выводы}
Период колебаний поршня в осцилляторе Фламмерсфельда получился $T = (335.5 \pm 1.2)\cdot10^{-3}\text{ с}$.
Показатель адиабаты воздуха получился $\gamma = (1.428\pm0.026)$, а количество эффективных степеней свободы
$i = (4.67\pm0.28)$, они близки к теоретическим, если счиать что воздух состоит преимущественно из двухатомных
молекул с замороженной колебательной степенью свободы ($\gamma = 1.4$, $i = 5$), но всё же не вкючают их в свои доверительные интервалы.
Это может быть связано с неучтённой погрешностью количества колебаний (например незарегистрированные половина периода).
Но несмотря на это полученные значения достаточно точны, их относительные погрешности не больше шести процентов.
\clearpage
\begin{figure}[H]
   \centering
   \includegraphics*[width=0.8\linewidth]{25274.jpg}
   \caption{Подписанные значения прямых измерений}
\end{figure}
\end{document}